\documentclass[10pt,a4paper]{amsart}
\usepackage[margin=19mm]{geometry}
\usepackage{amsmath,amssymb,mathtools}
\usepackage[scr=rsfs,cal=euler]{mathalfa}
\usepackage{xfrac}
\usepackage[unicode,hidelinks,bookmarks=false]{hyperref}
\newcommand{\ie}{i.e.\ }
\newcommand{\cf}{cf.\ }
\newcommand{\resp}{respectively\ }
\newcommand{\Subsection}{\subsection}
\providecommand{\nobreakdash}{\nobreak\hbox{-}}
\setlength{\emergencystretch}{2em}
\multlinegap0pt
\usepackage{bm}
\usepackage{cancel}
\usepackage{amssymb}
\usepackage{dsfont}
\usepackage{mathtools}
\usepackage{float}
\usepackage[shortlabels]{enumitem}
\usepackage{algorithm}\usepackage{algpseudocode}
\usepackage{xcolor}
\newtheorem{assumption}{Assumption}
\newtheorem{lemma}{Lemma}
\DeclareMathOperator{\E}{\mathbb{E}}
\DeclareMathOperator{\Q}{\mathbb{Q}}
\DeclareMathOperator{\R}{\mathbb{R}}
\def\mo{\mathfrak{o}}
\title{\vspace{-10mm} Cylindrical Projections of Occupied Diffusions \vspace{-3mm}}
\author{Valentin Tissot-Daguette, Xin Zhang}
\date{Source: arXiv:2604.25001v2 (CC BY 4.0)}
\begin{document}
\maketitle

\noindent\emph{Source and licence.} \vspace{-10mm} Cylindrical Projections of Occupied Diffusions \vspace{-3mm}. Version arXiv:2604.25001v2 (CC BY 4.0), distributed by arXiv under the Creative Commons Attribution 4.0 International licence, http://creativecommons.org/licenses/by/4.0/, as recorded in the arXiv licence metadata for this exact version and shown on the abstract page https://arxiv.org/abs/2604.25001v2. Full licence notice: \texttt{licenses/arxiv-2604.25001-CC-BY-4.0.txt}.\par\smallskip

\noindent\textit{Editorial scope.} Counted unit: the article's lemma lem:stopP (source0.tex lines 447--453). The article's complete original proof of it is delivered with it (source0.tex lines 455--457, one \texttt{proof} environment of 3 lines). No mathematics is written, repaired or re-derived: the statement and the proof are the article's own lines, in the article's own notation, and no retained line is substituted or re-flowed.  Retained uncounted as the source-local context the proof needs: lines 318--325; lines 1232--1232.  Nothing was omitted from any retained span, so the package carries no omission record.  No retained line contains a citation, so the wrapper carries no bibliography and no bibliography entry.  No retained line contains a figure environment, an image-inclusion command or drawing code, so no image is delivered and the package declares no asset dependency.  Editorial formatting only, in addition: an AMS \texttt{amsart} preamble that restates the article's own declarations and macros used by the retained lines, namely the environments \texttt{assumption}, \texttt{lemma} on separate counters, together with the standard packages those lines need.  The retained environments print in this document as Assumption 1, Lemma 1 in order of appearance, whereas the article prints the same items with the numbering of its own full document.  The complete native source of the article as distributed in the arXiv e-print for this exact version is bundled unchanged as \texttt{sources/199260/source0.tex}.
 \begin{assumption}\label{asm:growthLip}
     The coefficients $\lambda, b, \sigma$ satisfy the following linear growth and Lipschitz conditions: there exists $L> 0$ such that for all $\phi \in \{\lambda, b, \sigma\}$,
     \begin{align*}
        | \phi(\mo,x) | &\le L(1 + \lVert (\mo,x) \rVert ), \\[0.5em]
        | \phi(\mo,x) - \phi(\mo',x') | &\le L\lVert (\mo - \mo',x - x') \rVert.
     \end{align*}
     When $\phi = \sigma$, then $|\sigma(\mo,x)|$ denotes the Frobenius norm of $\sigma(\mo,x) \in \R^{d\times d}$. 
\end{assumption}

\subsection{Lemma~\ref{lem:stopP}} \label{app:Lemma}

\begin{lemma}\label{lem:stopP}
  Under Assumption~\ref{asm:growthLip}, for any $\gamma \in (0,1)$, $R>2$,
    \begin{align}\label{eq:lenglart}
       \Q( \tau_R \leq T)  \leq \frac{ 4 e^{\iota \gamma T} (2-\gamma) \E^{\Q}[(1+\lVert(\mo ,x)\rVert|)^{3\gamma}]}{R^{3 \gamma}(1-\gamma) },
    \end{align}
    where $\iota$ is a constant that only depends on $b,\sigma$. 
\end{lemma}

\begin{proof}
    See Appendix~\ref{app:Lemma}.
\end{proof}

\end{document}
